\scanpage{25}
\mspage{24}
% @ opening
Mais, \deleted{\unsure{on}} on arrive aux réciproques

% @ upper left, oblique
% slanted
\margin{Attention, Ces réciproques valent dans les exemples-types \unsure{unistrates|uni-strates} (théorie topologique, ou modérée) mais pas pour les réseaux p.ex.}

% @ Cont 5
\uline{\unsure{Cont|Conf}} 5 Soient $X,Y\in M$ deux multistrates, et considérons $\deleted{\unsure{L_X|\widehat L_X}}\widehat{\unsure{L|\mathcal{L}}}_X,\widehat L_Y\subset\unsure{\mathcal P|P}(L)$ les ens de parties de $L$ correspondantes. Ils permettent, grâce aux scholies de récupérer les strates ``ouvertes''

% @ a), first condition
a) Pour que $X\ll Y$, il f. et s. que \deleted{\unsure{tout} $\forall A\ \unsure{L}\ \forall A\in L_X$} toute strate ouverte de $X$ soit contenue dans une strate ouverte de $Y$

% @ left, beside (a) and (b)
% slanted
\margin{Il en résulte l'énoncé correspondant, pour deux figures $F,F'$, \unsure{déterminé|déterminés|déterminées} \unsure{par} $\widehat L_F,\widehat L_{F'}$.}

% @ b), second condition
b) Pour que $X$ et $Y$ soient compatibles, il f. et s. que pour toute strate ouverte de $X$ et toute strate ouverte de $Y$, \deleted{on ait} celles-ci soient disjointes ou identiques.

% @ Cor., two lines
\uline{Cor.}
\[
\begin{aligned}
\text{\deleted{b)} a)}\quad &\widehat L_X\unsure{\ne|=}\widehat L_Y\ \deleted{\unsure{\Longrightarrow|\text{donc}}}\Longleftrightarrow\ X=Y\\
\text{b)}\quad &\widehat L_X\subset\widehat L_Y\ \deleted{\unsure{\widehat L_X|\widehat L_X\,\widehat L_X}}\ \text{\unsure{est} fermé dans }\widehat L_Y\ \Longleftrightarrow\ X\leq Y
\end{aligned}
\]

% @ Scholie, bottom
\uline{Scholie} On récupère $M$, avec $\ll$, $\unsure{R|\mathcal R}$ (compatibilité), en termes de $L$, et de l'ens \inserted{$\widehat M$} des $\widehat L_X$
\[
\widehat M=\operatorname{Im}\bigl(\underbrace{M\xrightarrow{X\mapsto\widehat L_X}\mathcal P(\mathcal P}_{\text{application \uline{injective}}}(L))\ \in\mathcal P(\mathcal P(\mathcal P(L)))
\]
